%%%%%%%%%%%%%%%%%%%%%%%%%%%%%%%%%%%%%%%%%
% Medium Length Professional CV
% LaTeX Template
% Version 2.0 (8/5/13)
%
% This template has been downloaded from:
% http://www.LaTeXTemplates.com
%
% Original author:
% Trey Hunner (http://www.treyhunner.com/)
%
% Important note:
% This template requires the resume.cls file to be in the same directory as the
% .tex file. The resume.cls file provides the resume style used for structuring the
% document.
%
%%%%%%%%%%%%%%%%%%%%%%%%%%%%%%%%%%%%%%%%%

%----------------------------------------------------------------------------------------
%   PACKAGES AND OTHER DOCUMENT CONFIGURATIONS
%----------------------------------------------------------------------------------------

\documentclass{resume} % Use the custom resume.cls style
\usepackage[left=0.7in,top=0.4in,right=0.7in,bottom=0.5in]
{geometry} % Document margins
\usepackage{url}
\usepackage{bookman}
\usepackage[T1]{fontenc}
\usepackage{hyperref}

%----------------------------------------------------------------------------------------
%   DOCUMENT START
%----------------------------------------------------------------------------------------
\name{Abhijeet Lokhande} % Your name
\designation{AI Engineer: Multi-Agent Systems \& LLM Fine-Tuning}

\address{07480801382 \\ London, United Kingdom}
\address{\href{mailto:abhijeetlokhande1996@gmail.com}{abhijeetlokhande1996@gmail.com} \\ https://www.linkedin.com/in/ablds/}

\begin{document}
\title{Anon Resume}
%----------------------------------------------------------------------------------------
%   EDUCATION SECTION
%----------------------------------------------------------------------------------------
\begin{rSection}{Summary}
AI engineer specializing in multi-agent orchestration (LangGraph) and LLM fine-tuning for production systems. Architected autonomous agent systems that reduced manual work by 70\% in financial modeling and document intelligence. Expert in model adaptation (LoRA, PEFT, fine-tuning) with 50K+ HuggingFace downloads. Proven ability to design end-to-end AI pipelines from data ingestion through model optimization to production deployment, with obsessive focus on reliability and performance.
\end{rSection}

\begin{rSection}{Technical Skills}

\begin{tabular}{ @{} >{\bfseries}l @{\hspace{6ex}} l }
Agents & {LangGraph, LangChain, LlamaIndex, Anthropic API, FastAPI}\\
ML/LLM & {PyTorch, Fine-tuning (LoRA, PEFT), Ollama, RAG, Vector embeddings}\\
Languages & {Python, SQL}\\
Data/DevOps & {PostgreSQL+pgvector, MongoDB, Docker, Git, Apache Hamilton}\\
MLOps & {Production monitoring, Error handling, Deployment pipelines}\\
Cloud & {AWS, GCP}\\
\end{tabular}


\end{rSection}

%----------------------------------------------------------------------------------------
%   Personal Details
%----------------------------------------------------------------------------------------


\begin{rSection}{Profile Links}
\begin{tabular}{ @{} >{\bfseries}l @{\hspace{6ex}} l }
Linkedin & https://www.linkedin.com/in/ablds/  \\
Github & https://github.com/abhijeetscode \\
Hugging Face & https://huggingface.co/ab-ai
\end{tabular}

\end{rSection}


%----------------------------------------------------------------------------------------
%   WORK EXPERIENCE SECTION
%----------------------------------------------------------------------------------------
\begin{rSection}{Open Source Contributions}
\begin{itemize}
    \item Built autonomous ReAct-based coding agent using LangGraph and Anthropic API that generates installable, tested software packages from plain-language markdown specs, achieving 9/9 pass rate across polyglot benchmark (Python, Go, Rust, Java, JS). Demonstrates tool-calling orchestration and agent reliability patterns. \url{https://github.com/abhijeetscode/Spec-to-Code-Agent}
    \item Fine-tuned Large Language Models (GPT with LoRA and BERT) for Personally Identifiable Information (PII) detection, with published models achieving 30,000+ downloads on HuggingFace. \url{https://huggingface.co/ab-ai/PII-Model-Phi3-Mini}
    \item Developed CLIP-style dual-encoder architecture for natural language image search, demonstrating experience with multimodal AI systems. \url{https://github.com/abhijeetscode/text-to-image-search}

\end{itemize}
\end{rSection}

\begin{rSection}{Work Experience}
\begin{rSubsection}{Built AI (FinTech \& AI Company)}{April 2022 to Present}{AI Engineer}{London, UK}
\item Autonomous DCF modeling agent: Created and deployed agentic system that automates Discounted Cash Flow financial model generation from raw financial data. Reduces analyst modeling time by 70\%. Integrated with Built AI platform; includes validation frameworks, error handling, and HITL escalation for high-stakes financial decisions.
\item Multi-agent architecture and orchestration: Designed and deployed autonomous agent systems using LangGraph for document intelligence, query interpretation, and financial analysis. Built with production-grade reliability: validation frameworks, error handling, HITL escalation, and cloud deployment on AWS.
\item LLM fine-tuning and model optimization: Fine-tuned GPT and BERT models for production use cases (PII detection, domain adaptation) using PyTorch. Implemented advanced techniques including LoRA and supervised fine-tuning. Optimized inference latency while maintaining accuracy; deployed models achieving 50K+ downloads on HuggingFace.
\item Production RAG systems and inference: Engineered document intelligence platform combining embeddings, vector databases (PostgreSQL + pgvector), and LLMs. Achieved 90\%+ retrieval accuracy with <2s query latency. Eliminated 80\% of manual document review through optimized pipeline design.
\item MLOps and cloud infrastructure: Rearchitected financial modeling engine from naive iteration to graph-based approach (Apache Hamilton + vectorization), achieving 85\% latency improvement. Designed complex data pipelines on AWS cloud linking datasets across sources with geolocation enrichment and ML inference. Built production-grade monitoring, error handling, and deployment pipelines.
\item Proactive automation and ROI: Discovered automation bottlenecks in analyst workflows; designed end-to-end solutions; measured and communicated business impact to leadership. Each initiative reduced manual work by 30-70\%.
\end{rSubsection}


\begin{rSubsection}{King's College London}{March 2021 to January 2022}{Research Associate}{London, UK}
\item Graph-based deep learning and model architecture: Developed novel models combining Graph Convolution Networks and Generative Adversarial Networks for complex optimization problems, achieving 63\% improvement in solution novelty. Demonstrates expertise in transformer-based and graph architectures applicable to structured data.
\item Data pipeline engineering and preprocessing: Designed and implemented robust preprocessing pipeline for clinical data at scale (1M+ rows), ensuring data integrity and completeness for downstream ML systems.
\item Web scraping and data collection: Utilized Scrapy and Beautiful Soup for large-scale data acquisition and processing.

\end{rSubsection}

\end{rSection}


%----------------------------------------------------------------------------------------
%   Personal Project
%----------------------------------------------------------------------------------------


\iffalse
\begin{rSection}{Relevant Project}
\begin{description}
  \item[$\cdot$] Solved infinite monkey problem using evolutionary algorithms.
  \item[$\cdot$] Automated a Pacman on a small grid using reinforcement learning.
  \item[$\cdot$] I built VAE to generate new images from the fashion MNIST dataset.
  \item[$\cdot$] I built a movie recommender system using cosine similarity.
\end{description}
\end{rSection}
\fi


%----------------------------------------------------------------------------------------
%   Achievements
%----------------------------------------------------------------------------------------
\iffalse
\begin{rSection}{Achievements}
\begin{description}
  \item[$\cdot$] Completed all level of Google FooBar.
  \item[$\cdot$] Got a certificate of appreciation from Merck \& Co.
  \item[$\cdot$] Arctic Code Vault Contributor of Github.
\end{description}
\end{rSection}
\fi
%----------------------------------------------------------------------------------------
%   Technical Experience
%----------------------------------------------------------------------------------------



%----------------------------------------------------------------------------------------
%   TECHNICAL STRENGTHS SECTION
%----------------------------------------------------------------------------------------
\begin{rSection}{Education}
% September 2020 - January 2022
\begin{rSubsection}{King's College London}{}{MSc in Data Science }{London, UK}
%\vspace{-.1cm}
%Overall GPA: 78\%
\vspace{-.65cm}
\item[]
\end{rSubsection}
\begin{rSubsection}{C-DAC: Centre for Development of Advanced Computing, India}{}{Post Graduate Diploma in Advanced Computing (PG-DAC)}{Pune, India}
%\vspace{-.1cm}
%Overall GPA: 78\%
\vspace{-.65cm}
\item[]
\end{rSubsection}
% June 2013 - June 2017
\begin{rSubsection}{University of Pune}{}{Bachelor of Engineering in Computer Science}{Pune, India}
%\vspace{-.1cm}
%Overall GPA: 78\%
\vspace{-.65cm}
\item[]
\end{rSubsection}
\end{rSection}

%----------------------------------------------------------------------------------------

\end{document}
